% ECONOMICS_PROBLEM: {"id": "D63-473", "title": "Maximal EF1 allocations for three agents with conflicting goods", "jel_code": "D63", "source_id": "OP-0197"}
% !TeX program = xelatex
\documentclass[11pt,a4paper]{article}
\usepackage[margin=23mm]{geometry}
\usepackage{fontspec}
\setmainfont{Latin Modern Roman}
\usepackage{amsmath,amssymb,mathtools}
\usepackage{xcolor,enumitem,booktabs,longtable,array,xurl,hyperref}
\definecolor{muted}{HTML}{53636C}
\hypersetup{colorlinks=true,urlcolor=blue,linkcolor=blue}
\setlength{\parindent}{0pt}
\setlength{\parskip}{5pt}
\newcommand{\R}{\mathbb R}
\newcommand{\E}{\mathbb E}
\newcommand{\N}{\mathbb N}
\newcommand{\ind}{\mathbf 1}
\newcommand{\Var}{\operatorname{Var}}
\newcommand{\Cov}{\operatorname{Cov}}
\newcommand{\supp}{\operatorname{supp}}
\DeclareMathOperator*{\argmin}{arg\,min}
\DeclareMathOperator*{\argmax}{arg\,max}
\newcommand{\entrymeta}[3]{{\sffamily\small\color{muted}\textbf{\texttt{#1}}\quad|\quad #2\par #3\par}}
\newcommand{\Problemref}[1]{\texttt{#1}}
\newcommand{\JELref}[2]{\texttt{#2} (JEL \texttt{#1})}
\begin{document}
\section*{Maximal EF1 allocations for three agents with conflicting goods}
\textbf{Problem ID:} \texttt{D63-473}\quad \textbf{JEL:} \texttt{D63}\par
% BEGIN ECONOMICS STATEMENT
\label{problem:OP-0197}
% GAME_THEORY_PROBLEM: {"local_id": "GT-067", "original_number": 67, "kind": "P", "entry_kind": "statement_only", "published": false, "review_required": true, "category": "game_theory"}
\label{game:GT-067}
{\sffamily\small\color{muted}
\noindent\textbf{ID: GT-067.}\quad\textbf{Category:} Constrained fair division.
\quad\textbf{Kind: P.}\quad\textbf{Status:} Open in the cited source.
\par}
\subsection*{Definitions and mathematical statement}
Let \(N=\{1,2,3\}\), let \(M\) be a finite set of goods, and let
\(H=(M,E)\) be a simple undirected graph. For each \(i\), let
\(v_i(S)=\sum_{g\in S}a_{ig}\), where \(a_{ig}\geq 0\).
A feasible partial allocation is a triple of pairwise disjoint sets
\(A_i\subseteq M\), each an independent set of \(H\).
Write \(U=M\setminus\bigcup_i A_i\). Such an allocation is
\emph{maximal} if
\[
 \forall g\in U\ \forall i\in N\ \exists h\in A_i:
       \{g,h\}\in E.
\]
Equivalently, no unallocated good can be added to any bundle while
preserving feasibility. It is EF1 if, for every \(i,j\in N\),
\[
 v_i(A_i)\geq v_i(A_j)
 \quad\text{or}\quad
 \exists g\in A_j:\ v_i(A_i)\geq v_i(A_j\setminus\{g\}).
\]

\paragraph{Question.}
Does every such instance admit a feasible partial allocation that is
both maximal and EF1?

\subsection*{Short English statement}
Can three agents divide conflicting goods fairly until no remaining
good can be assigned?

\subsection*{Variants and scope boundaries}
Maximal means that no single addition is feasible; it does not mean
maximum assigned cardinality or maximum welfare.
Completeness is not required, since a conflict graph need not admit a
proper three-coloring. EF1 compares the actual bundles using their
additive values.

\subsection*{Source relationship and status}
[S1, Section 4.1] leaves the three-agent additive case open.
Its impossibility results with more agents or different valuation
classes do not settle this restriction.

\subsection*{References}
\begingroup\footnotesize\raggedright
\noindent[S1] Ayumi Igarashi, Pasin Manurangsi, and Hirotaka Yoneda,
\emph{Dividing Conflicting Items Fairly}, IJCAI 2025;
arXiv:2506.14149v1.
\url{https://arxiv.org/html/2506.14149v1}.
Locators: Section 2 (feasibility, maximality, and EF1);
Section 4.1, paragraph following Theorem 11; Proposition 12.
\par\endgroup

\subsection*{Common conventions: Source mathematical conventions}

\textbf{Finite games.} For a finite set $A$, $\Delta(A)$ is its probability
simplex. Players choose independent mixed strategies in Nash equilibrium;
correlation is allowed only when explicitly part of the solution concept.
In a game with payoff functions $u_i$ and strategy profile $x$, an additive
$\varepsilon$ Nash equilibrium satisfies
\[
 u_i(x)\geq u_i(x_i',x_{-i})-\varepsilon
 \quad\text{for every player $i$ and every admissible deviation $x_i'$}.
\]
For numerical approximation, payoffs are normalized as locally specified.
An $\varepsilon$ payoff residual is distinct from distance at most
$\varepsilon$ to the set of exact equilibria. Point convergence, convergence
of empirical play, and convergence in distance to an equilibrium set are
also distinct.

\textbf{Computational representations.} Unless stated otherwise, a complexity
question takes rational data with binary encoding; polynomial time is measured
in total bit length. A fixed-accuracy polynomial algorithm is not necessarily
polynomial in $1/\varepsilon$, and neither is necessarily polynomial in
$\log(1/\varepsilon)$. Succinct games and value, demand, or sampling oracles
must declare their interfaces. Exact real solutions require an explicit
representation; an unspecified list of arbitrary real numbers is not a
finite computational output. A claimed algorithmic barrier is meaningful
only for its specified model and reductions.

\textbf{Dynamic games.} Histories, information, observation, and allowable
randomization are part of the model. A stationary strategy depends only on
the current state; a positional strategy in a deterministic graph game
chooses one action at each controlled vertex. Nash equilibrium at an initial
state and equilibrium after every history are different requirements.
An expectation of a pathwise $\liminf$ payoff, a $\liminf$ of expected averages,
a limiting expected average, and a uniform finite-horizon equilibrium are
different criteria. None is silently substituted for another.

\textbf{Allocation and mechanisms.} A complete allocation assigns every item
exactly once unless otherwise stated. Zero-valued goods, negative values,
capacity constraints, feasibility, lotteries, and copies of goods affect
fairness definitions and are specified locally. Dominant-strategy incentive
compatibility, Bayesian incentive compatibility, universal truthfulness,
and truthfulness in expectation impose different inequalities. Whenever
an objective ratio has zero denominator, its convention is stated or those
instances are excluded.

\textbf{Infinite-dimensional models.} Expectations, integrals, stochastic
equations, and optimization objectives require their local regularity,
integrability, filtrations, and admissible domains. A backward--forward
mean-field system is a boundary-value problem; it is not an initial-value
semiflow without an additional construction. Long-time, large-population,
and vanishing-noise limits require an order of limits and a convergence
topology. If a minimum need not be attained, the objective is its infimum
or supremum and the approximation requirement is stated separately.
% END ECONOMICS STATEMENT
\end{document}
